\section{Original Theoretical Contribution}
\label{ch:original_theoretical_contribution}

\subsection{Overview of the contribution}

Chapter~2 identified a gap in the existing perturbation-learning literature: although NP is often expected to be more efficient than WP, the relationship between the two methods has not been fully characterized at the estimator level in nonlinear multilayer networks. This chapter addresses that gap by developing a direct theoretical comparison between weight-level and node-level perturbations. The goal is not only to compare WP and NP more precisely, but also to use this comparison to motivate a theoretically grounded variant of node perturbation.

The starting point for this analysis is the local reparameterization trick of \textcite{kingma2015variational}, although their setting is different from ours. They study stochastic variational inference in Bayesian neural networks, where the goal is to learn a distribution over weights rather than a single deterministic set of weights. The parameters of this distribution are optimized using backpropagation, so the learning rule itself is not perturbation-based in the sense considered in this thesis. However, the variational objective contains an expectation over random weights, and this expectation is estimated using Monte Carlo samples from the weight distribution. In this respect, their estimator has an important similarity to WP: both involve sampling noise at the weight level and using the resulting stochastic forward pass to estimate a gradient. Kingma et al. show that, in their setting, this weight-level sampling can be replaced by sampling the equivalent activation-level noise directly, and that doing so reduces the variance of the gradient estimator. This motivates the question studied here: whether the same line of reasoning can be used to derive a node-level perturbation rule that matches the noise induced by WP, and then compare the resulting estimator to WP directly.

This chapter applies that idea in the perturbation-learning setting. We first derive the distribution of the pre-activation noise induced by Gaussian WP in a nonlinear multilayer network. This shows that WP does not induce fixed-variance node noise; instead, the variance depends on the norm of the incoming activation vector. We then use this distribution to define input-scaled node perturbation (IS-NP), a node perturbation rule that samples directly from the node-level noise distribution induced by WP. Finally, we show that WP and IS-NP estimate the same perturbation-smoothed objective in expectation, before proving that IS-NP has variance smaller than or equal to WP for each coordinate. Together, these results provide a direct theoretical comparison between weight-level and node-level perturbations, and identify IS-NP as a variance-reduced node-level version of WP.
\subsection{Pre-activation noise induced by WP}

To compare WP and NP more directly, we first express the effect of WP at the same level where NP applies perturbations: the pre-activations. This allows us to characterize how independent weight perturbations propagate into neuron-level noise, and we can then compare this induced noise with the perturbations injected directly by NP.
Recall from Section~\ref{sec:weight_perturbation} that we modeled perturbations as a deterministic variable plus a scale times unit Gaussian noise. For neural networks, this corresponds to perturbing each weight independently as
\[
\mathbf{W}^{\ell\prime}
=
\mathbf{W}^\ell + \sigma \boldsymbol{\xi}_W^\ell,
\qquad
\xi_{W,ij}^\ell \sim \mathcal{N}(0,1)\ \text{i.i.d.},
\]
where \(\xi_{W,ij}^\ell\) denotes the \((i,j)\)-entry of the unit-noise matrix \(\boldsymbol{\xi}_W^\ell\).

The pre-activation at neuron $j$ in layer $\ell$ is given by
\[
z^\ell_j = \sum_i W^\ell_{ij} x^{\ell-1}_i.
\]
After perturbation, this becomes

\[
z^{\ell\prime}_j =
\sum_i W^{\ell\prime}_{ij} x^{\ell-1\prime}_i
=
\sum_i (W^\ell_{ij} + \sigma \xi_{W,ij}^\ell) \, x^{\ell-1\prime}_i.
\]
This expression separates into a deterministic component and a stochastic component induced by the weight perturbations:
\[
z^{\ell\prime}_j
=
\sum_i W^\ell_{ij} x^{\ell-1\prime}_i
+
\sigma \sum_i \xi_{W,ij}^\ell x^{\ell-1\prime}_i.
\]

Define the induced pre-activation noise as
\begin{equation}
\Delta z^\ell_j
:=
\sigma \sum_i \xi_{W,ij}^\ell x^{\ell-1\prime}_i.
\label{eq:preact_noise}
\end{equation}

We can now write the full pre-activation in compact form as
\[
z^{\ell\prime}_j = (\mathbf{W}^\ell \mathbf{x}^{\ell-1\prime})_j + \Delta z^\ell_j,
\]

\(\Delta z^\ell_j\) is the effect of weight perturbation at the pre-activation level that we wanted to express. We now want to characterize the distribution of this noise term.

Condition on the incoming perturbed activations \(\mathbf{x}^{\ell-1\prime}\). Under this conditioning, the
coefficients $x^{\ell-1\prime}_i$ are fixed, while the weight-noise variables
$\xi_{W,ij}^\ell$ remain independent standard Gaussians. Therefore each term
$\sigma x^{\ell-1\prime}_i \xi_{W,ij}^\ell$ is Gaussian:
\begin{equation}
\sigma x^{\ell-1\prime}_i \xi_{W,ij}^\ell
\mid \mathbf{x}^{\ell-1\prime}
\sim
\mathcal{N}\!\left(0, \sigma^2 (x^{\ell-1\prime}_i)^2\right).
\end{equation}
Since $\Delta z^\ell_j$ is a sum of independent Gaussian random variables, it is also Gaussian.
Its conditional mean is
\begin{equation}
\mathbb{E}\!\left[\Delta z^\ell_j \mid \mathbf{x}^{\ell-1\prime}\right]
=
\sigma \sum_i x^{\ell-1\prime}_i
\mathbb{E}\!\left[\xi_{W,ij}^\ell\right]
=
0,
\end{equation}
and its conditional variance is
\begin{equation}
\operatorname{Var}\!\left(\Delta z^\ell_j \mid \mathbf{x}^{\ell-1\prime}\right)
=
\sigma^2 \sum_i (x^{\ell-1\prime}_i)^2
\operatorname{Var}\!\left(\xi_{W,ij}^\ell\right)
=
\sigma^2 \left\|\mathbf{x}^{\ell-1\prime}\right\|^2 .
\end{equation}
Thus, the pre-activation noise induced by Gaussian weight perturbation satisfies
\begin{equation}
\Delta z^\ell_j \mid \mathbf{x}^{\ell-1\prime}
\sim
\mathcal{N}\!\left(0, \sigma^2 \left\|\mathbf{x}^{\ell-1\prime}\right\|^2\right).
\end{equation}

This shows that independent Gaussian perturbations at the weight level do not produce arbitrary pre-activation noise. Instead, their combined effect at each neuron is itself Gaussian, with a variance determined by the norm of the incoming activation vector. In other words, WP induces structured node-level noise: neurons receiving larger incoming activations are perturbed more strongly, while neurons receiving smaller incoming activations are perturbed less. This induced distribution is the key object we will use to compare WP with node perturbation directly.

\subsection{Input-scaled node perturbation}

The result above suggests a natural modification of standard node perturbation. Recall from Section~2.5 that standard NP injects Gaussian noise with the same perturbation scale at every node. That is, although the perturbation scale is a hyperparameter, it is fixed across neurons and does not depend on the magnitude of the incoming activation vector. This means that standard NP does not reproduce the pre-activation noise induced by WP: as shown in Section~3.2, Gaussian weight perturbations induce node-level noise whose variance scales with $\|\mathbf{x}^{\ell-1\prime}\|^2$.

One way to approximate this input-dependent scaling is fan-in-scaled NP. If the incoming activations are independent standard Gaussians, $x_i^{\ell-1} \sim \mathcal{N}(0,1)$, then
\begin{equation}
\mathbb{E}\!\left[\left\|\mathbf{x}^{\ell-1}\right\|^2\right]
=
\mathbb{E}\!\left[\sum_{i=1}^{N_{\ell-1}} (x_i^{\ell-1})^2\right]
=
N_{\ell-1}.
\end{equation}
Under this assumption, the induced WP noise has expected variance $\sigma^2 N_{\ell-1}$ at each node. This is exactly the scaling argument used by \textcite{werfel2005learning}, who choose the NP noise variance to be a factor $N$ larger than the WP noise variance in order to make the two stochastic variants directly comparable. Fan-in-scaled NP therefore matches the WP-induced node noise under this distributional assumption, but only at the level of the expected squared input norm.

We now remove this expectation-level approximation by matching the WP-induced pre-activation noise directly. This leads to the method proposed in this thesis, which we call input-scaled node perturbation (IS-NP). To the best of our knowledge, this specific perturbation rule has not been derived previously; in particular, it does not appear among the perturbation-learning methods reviewed in Table~\ref{tab:perturbation_update_rules}. For each neuron $j$ in layer $\ell$, we sample
\begin{equation}
\Delta z^\ell_j
=
\sigma \left\|\mathbf{x}^{\ell-1\prime}\right\| \xi^\ell_{z,j},
\qquad
\xi^\ell_{z,j} \sim \mathcal{N}(0,1).
\end{equation}
This gives the same conditional distribution as the WP-induced pre-activation noise derived in Section~3.2:
\begin{equation}
\Delta z^\ell_j \mid \mathbf{x}^{\ell-1\prime}
\sim
\mathcal{N}\!\left(0, \sigma^2 \left\|\mathbf{x}^{\ell-1\prime}\right\|^2\right).
\end{equation}

We now derive the corresponding update rule. Since the perturbation standard deviation is
$\sigma \|\mathbf{x}^{\ell-1\prime}\|$, the score-function term for the pre-activation perturbation is
\begin{equation}
\frac{\Delta z^\ell_j}{\sigma^2 \left\|\mathbf{x}^{\ell-1\prime}\right\|^2}
=
\frac{\xi^\ell_{z,j}}{\sigma \left\|\mathbf{x}^{\ell-1\prime}\right\|}.
\end{equation}
Let $\Delta \ell = L_{\mathrm{pert}} - b$ denote the scalar learning signal. The resulting gradient estimate with respect to the pre-activation is therefore
\begin{equation}
\hat{g}_{z^\ell_j}
=
\Delta \ell
\frac{\xi^\ell_{z,j}}{\sigma \left\|\mathbf{x}^{\ell-1\prime}\right\|}.
\end{equation}
Using the local derivative
\begin{equation}
\frac{\partial z^\ell_j}{\partial W^\ell_{ij}}
=
x^{\ell-1\prime}_i,
\end{equation}
we obtain the corresponding weight-gradient estimator
\begin{equation}
\hat{g}^{\mathrm{IS\text{-}NP}}_{ij,\ell}
=
\Delta \ell
\frac{\xi^\ell_{z,j} x^{\ell-1\prime}_i}
{\sigma \left\|\mathbf{x}^{\ell-1\prime}\right\|}.
\label{eq:isnp_estimator}
\end{equation}
Applying gradient descent gives the update rule
\begin{equation}
\Delta W^\ell_{ij}
=
-\eta \Delta \ell
\frac{\xi^\ell_{z,j} x^{\ell-1\prime}_i}
{\sigma \left\|\mathbf{x}^{\ell-1\prime}\right\|}.
\end{equation}

Thus, unlike standard NP, the perturbation scale is not fixed across neurons. Unlike fan-in-scaled NP, it is also not determined only by the number of incoming units. Instead, IS-NP scales the perturbation by the actual norm of the incoming activation vector for each layer and input example. This does not by itself imply that IS-NP is preferable to other NP variants in practice, but it makes IS-NP directly comparable to WP because it samples the same pre-activation noise distribution induced by Gaussian weight perturbations.
\subsection{Expected equivalence between WP and IS-NP}

Let $F_{\mathrm{WP}}(W)$ denote the perturbation-smoothed objective optimized by WP:
\begin{equation}
F_{\mathrm{WP}}(W)
=
\mathbb{E}_{\Delta W}
\left[
L\!\left(\mathrm{forward}(W+\Delta W)\right)
\right].
\end{equation}
Although WP samples perturbations in weight space, the loss only depends on these perturbations through the perturbed forward pass. Equivalently, we can describe the same objective in terms of the induced pre-activation perturbations. Let $p_{\mathrm{WP}}(\Delta z \mid x)$ denote the distribution over pre-activation perturbations induced by WP. Then
\begin{equation}
F_{\mathrm{WP}}(W)
=
\mathbb{E}_{\Delta z \sim p_{\mathrm{WP}}(\Delta z \mid x)}
\left[
L\!\left(\mathrm{forward}(W;\Delta z)\right)
\right].
\end{equation}

By construction, IS-NP samples directly from this same induced distribution:
\begin{equation}
p_{\mathrm{IS\text{-}NP}}(\Delta z \mid x)
=
p_{\mathrm{WP}}(\Delta z \mid x).
\end{equation}
Therefore,
\begin{align}
F_{\mathrm{IS\text{-}NP}}(W)
&=
\mathbb{E}_{\Delta z \sim p_{\mathrm{IS\text{-}NP}}(\Delta z \mid x)}
\left[
L\!\left(\mathrm{forward}(W;\Delta z)\right)
\right] \\
&=
\mathbb{E}_{\Delta z \sim p_{\mathrm{WP}}(\Delta z \mid x)}
\left[
L\!\left(\mathrm{forward}(W;\Delta z)\right)
\right] \\
&=
F_{\mathrm{WP}}(W).
\end{align}
Thus, WP and IS-NP define the same perturbation-smoothed objective, even though they generate the perturbations in different spaces.
Since both methods use unbiased score-function estimators of this same objective, their expected gradient estimates are equal:
\begin{equation}
\mathbb{E}\!\left[\hat{g}^{\mathrm{WP}}\right]
=
\nabla F(W)
=
\mathbb{E}\!\left[\hat{g}^{\mathrm{IS\text{-}NP}}\right].
\end{equation}
In other words, WP and IS-NP have the same expected update direction, because they are both unbiased estimators of the gradient of the same objective function. The same equivalence does not generally hold for standard NP or fan-in-scaled NP. While these methods also produce unbiased gradient estimators of perturbation-smoothed objective functions, the functions themselves can be slightly different because they induce different distributions over perturbed forward passes than WP. As $\sigma \to 0$, all these smoothed objectives approach the original deterministic loss, so this difference vanishes in the infinitesimal-noise limit. For nonzero $\sigma$, however, the objectives are not exactly identical: the same perturbed forward pass may be assigned different probability under different perturbation schemes. Therefore, comparing WP with standard NP or fan-in-scaled NP is not a pure estimator-level comparison, since the estimators may target different functions. IS-NP removes this ambiguity by matching the WP-induced perturbation distribution directly. Since WP and IS-NP have the same expectation, the natural remaining question is how much the two estimators vary around this shared expected direction. This gives IS-NP a useful theoretical role beyond being a new practical variant: it provides a node-level estimator that can be compared with WP without changing the underlying objective being estimated. The variance comparison in the next section is therefore a clean estimator-level comparison, rather than a comparison between methods that may optimize slightly different smoothed objectives.
\subsection{Variance reduction theorem}

The previous section showed that IS-NP and WP are matched in expectation. We now show that they differ in variance. More specifically, IS-NP can be viewed as the conditional expectation of WP given the induced pre-activation noise field. This means that IS-NP removes weight-noise variability that affects the WP update but does not affect the actual perturbed forward pass. The result is a direct variance reduction by the law of total variance.

For a single coordinate \(W^\ell_{ij}\), the WP estimator can be written as
\begin{equation}
\widehat{g}^{\mathrm{WP}}_{ij,\ell}
=
\Delta \ell
\frac{\xi^\ell_{W,ij}}{\sigma}.
\label{eq:wp_estimator_component}
\end{equation}

\begin{proposition}
For Gaussian weight perturbations and the corresponding input-scaled node perturbation estimator defined in Eq.~\eqref{eq:isnp_estimator}, the coordinate-wise variance of IS-NP is no larger than that of WP:
\[
\mathrm{Var}\!\left(\widehat{g}^{\mathrm{IS\text{-}NP}}_{ij,\ell}\right)
\le
\mathrm{Var}\!\left(\widehat{g}^{\mathrm{WP}}_{ij,\ell}\right).
\]
\end{proposition}

\begin{proof}
We consider the weight perturbation estimator in \eqref{eq:wp_estimator_component}, and analyze its variance by conditioning on the induced pre-activation noise. Recall from \eqref{eq:preact_noise} that $\Delta z^\ell_j$ is a linear combination of the weight perturbations.

Condition on the full induced noise field \(\boldsymbol{\Delta Z}=(\boldsymbol{\Delta z}^1,\dots,\boldsymbol{\Delta z}^L)\). Since the perturbed forward pass satisfies
\[
\mathbf{z}^{\ell\prime}
=
\mathbf{W}^\ell \mathbf{x}^{\ell-1\prime}
+
\boldsymbol{\Delta z}^\ell,
\]
the entire perturbed trajectory \((\mathbf{x}^{1\prime},\dots,\mathbf{x}^{L\prime})\) is uniquely determined by \(\boldsymbol{\Delta Z}\). Hence both the scalar learning signal \(\Delta \ell(\boldsymbol{\Delta Z})\) and all activations \(\mathbf{x}^{\ell\prime}\) are deterministic given \(\boldsymbol{\Delta Z}\).

Now fix \((i,j,\ell)\). Given $\boldsymbol{\Delta Z}$, the incoming activation vector $\mathbf{x}^{\ell-1\prime}$ is fixed, and the only component of the induced noise field that contains information about $\xi^\ell_{W,ij}$ is $\Delta z^\ell_j$. From \eqref{eq:preact_noise}, $\Delta z^\ell_j$ is a linear combination of Gaussian weight perturbations. Hence, $\xi^\ell_{W,ij}$ and $\Delta z^\ell_j$ are jointly Gaussian. Since both variables have mean zero, the conditional expectation is linear and given by
\begin{equation}
\mathbb{E}\!\left[\xi^\ell_{W,ij} \mid \boldsymbol{\Delta Z}\right]
=
\frac{
\operatorname{Cov}\!\left(\xi^\ell_{W,ij}, \Delta z^\ell_j\right)
}{
\operatorname{Var}\!\left(\Delta z^\ell_j\right)
}
\Delta z^\ell_j .
\end{equation}
The covariance term is
\begin{align}
\operatorname{Cov}\!\left(\xi^\ell_{W,ij}, \Delta z^\ell_j\right)
&=
\operatorname{Cov}\!\left(
\xi^\ell_{W,ij},
\sigma \sum_k \xi^\ell_{W,kj} x_k^{\ell-1\prime}
\right) \\
&=
\sigma \sum_k x_k^{\ell-1\prime}
\operatorname{Cov}\!\left(\xi^\ell_{W,ij}, \xi^\ell_{W,kj}\right) \\
&=
\sigma x_i^{\ell-1\prime},
\end{align}
where the last equality follows from the independence and unit variance of the Gaussian weight perturbations. Similarly,
\begin{equation}
\operatorname{Var}\!\left(\Delta z^\ell_j\right)
=
\sigma^2 \left\|\mathbf{x}^{\ell-1\prime}\right\|^2.
\end{equation}
Substituting these expressions gives
\begin{equation}
\mathbb{E}\!\left[\xi^\ell_{W,ij} \mid \boldsymbol{\Delta Z}\right]
=
\frac{x_i^{\ell-1\prime}}{\sigma \left\|\mathbf{x}^{\ell-1\prime}\right\|^2}
\Delta z^\ell_j.
\end{equation}
Therefore,
\begin{equation}
\mathbb{E}\!\left[
\widehat{g}^{\mathrm{WP}}_{ij,\ell}
\mid
\boldsymbol{\Delta Z}
\right]
=
\frac{\Delta \ell(\boldsymbol{\Delta Z})}{\sigma}
\mathbb{E}\!\left[\xi^\ell_{W,ij} \mid \boldsymbol{\Delta Z}\right]
=
\Delta \ell(\boldsymbol{\Delta Z})
\frac{x_i^{\ell-1\prime}\Delta z^\ell_j}
{\sigma^2 \left\|\mathbf{x}^{\ell-1\prime}\right\|^2}.
\end{equation}

Recall from \eqref{eq:isnp_estimator} that the IS-NP estimator is
\begin{equation}
\widehat{g}^{\mathrm{IS\text{-}NP}}_{ij,\ell}
=
\Delta \ell
\frac{\xi^\ell_{z,j} x_i^{\ell-1\prime}}
{\sigma \left\|\mathbf{x}^{\ell-1\prime}\right\|}.
\end{equation}
Since the IS-NP perturbation satisfies
\begin{equation}
\Delta z^\ell_j
=
\sigma \left\|\mathbf{x}^{\ell-1\prime}\right\|\xi^\ell_{z,j},
\end{equation}
the same estimator can equivalently be written as
\begin{equation}
\widehat{g}^{\mathrm{IS\text{-}NP}}_{ij,\ell}
=
\Delta \ell
\frac{x_i^{\ell-1\prime}\Delta z^\ell_j}
{\sigma^2 \left\|\mathbf{x}^{\ell-1\prime}\right\|^2}.
\end{equation}
Hence,
\begin{equation}
\mathbb{E}\!\left[
\widehat{g}^{\mathrm{WP}}_{ij,\ell}
\mid
\boldsymbol{\Delta Z}
\right]
=
\widehat{g}^{\mathrm{IS\text{-}NP}}_{ij,\ell}.
\end{equation}

The law of total variance states that for random variables \(Y\) and \(X\),
\[
\operatorname{Var}(Y)
=
\operatorname{Var}\!\left(\mathbb{E}[Y\mid X]\right)
+
\mathbb{E}\!\left[\operatorname{Var}(Y\mid X)\right].
\]
Applying this identity with \(Y=\widehat{g}^{\mathrm{WP}}_{ij,\ell}\) and \(X=\boldsymbol{\Delta Z}\),
\[
\mathrm{Var}(\widehat{g}^{\mathrm{WP}}_{ij,\ell})
=
\mathrm{Var}\!\bigl(\mathbb{E}[\widehat{g}^{\mathrm{WP}}_{ij,\ell}\mid \boldsymbol{\Delta Z}]\bigr)
+
\mathbb{E}\!\bigl[\mathrm{Var}(\widehat{g}^{\mathrm{WP}}_{ij,\ell}\mid \boldsymbol{\Delta Z})\bigr].
\]
Using the identity above,
\[
\mathrm{Var}(\widehat{g}^{\mathrm{WP}}_{ij,\ell})
=
\mathrm{Var}(\widehat{g}^{\mathrm{IS\text{-}NP}}_{ij,\ell})
+
\mathbb{E}\!\bigl[\mathrm{Var}(\widehat{g}^{\mathrm{WP}}_{ij,\ell}\mid \boldsymbol{\Delta Z})\bigr].
\]
Since the second term is nonnegative, it follows that
\[
\mathrm{Var}(\widehat{g}^{\mathrm{IS\text{-}NP}}_{ij,\ell})
\le
\mathrm{Var}(\widehat{g}^{\mathrm{WP}}_{ij,\ell}).
\qedhere
\]
\end{proof}

The intuition behind this result is simple. Once the induced pre-activation noise field $\boldsymbol{\Delta Z}$ is fixed, the perturbed forward pass is fixed. This means that the activations, the loss difference $\Delta \ell$, and the IS-NP update are all determined. In WP, however, the update still depends on the individual weight perturbations that produced the same $\boldsymbol{\Delta Z}$. These remaining weight-level perturbations no longer affect the forward pass or the loss signal; they only change the WP update itself. They therefore add variability that does not help learning. From the perspective of estimating the same learning signal, this variability is simply extra noise.

This variance-reduction argument is closely related to the local reparameterization trick of \textcite{kingma2015variational}. Their setting is different: they study stochastic variational inference in Bayesian neural networks, where uncertainty over global weight parameters is translated into local activation noise to obtain lower-variance gradient estimates. The underlying principle is nevertheless the same. In both cases, weight-level noise induces a lower-dimensional distribution over activations or pre-activations, and sampling this local noise directly preserves the relevant distribution over perturbed forward passes while removing weight-level randomness that does not affect the forward computation. The IS-NP construction can therefore be viewed as applying the same variance-reduction idea to perturbation-based learning, where it provides a direct estimator-level comparison between WP and a node-level perturbation rule.

This result addresses the knowledge gap identified in Section~2.7 and fulfils the second research objective, namely to further develop the theory of perturbation-based learning. The literature reviewed earlier provides intuition for why node perturbation may be more efficient than weight perturbation, but it does not give a general estimator-level comparison between the two in nonlinear multilayer networks. The analysis in this chapter provides such a comparison by deriving IS-NP from the pre-activation noise induced by WP, showing that WP and IS-NP estimate the same perturbation-smoothed objective in expectation, and proving that IS-NP has variance smaller than or equal to WP for each coordinate. Thus, the relationship between weight-level and node-level perturbations has been clarified theoretically, and this understanding has been used to develop a theoretically motivated variant of node perturbation. Since IS-NP estimates the same perturbation-smoothed objective as WP while having variance smaller than or equal to WP for each coordinate, the result provides a theoretical reason to prefer IS-NP over WP for nonlinear multilayer networks in general.

\subsection{Why IS-NP may improve over existing NP variants}
The variance theorem above compares IS-NP directly with WP, but it does not establish how IS-NP compares to other node perturbation variants. Standard NP and fan-in-scaled NP use different perturbation distributions, and therefore generally correspond to perturbation-smoothed objectives that are similar, but not identical, to those of WP and IS-NP. For this reason, we do not derive an analytical variance comparison between IS-NP and these methods. Instead, we use the scaling arguments developed above to reason about which perturbation rule is likely to be better behaved in practice, before testing this empirically in Chapter~5.

Standard NP uses the same perturbation scale at every node, regardless of the incoming activation norm. This can lead to poorly scaled updates. If $\|\mathbf{x}^{\ell-1}\|$ is large, a fixed pre-activation perturbation may have only a small relative effect on the network output, while the resulting weight update can still become large because it is multiplied by the incoming activation. Conversely, if $\|\mathbf{x}^{\ell-1}\|$ is small, the same fixed perturbation may create a large relative change in the pre-activation, while the corresponding weight update is small. In both cases, the perturbation scale is disconnected from the activation scale that determines the weight update. This mismatch is expected to be most problematic in networks where activation norms vary substantially across layers, samples, or training time, for example when layer widths differ or activations are not explicitly normalized. In such cases, a fixed perturbation scale will sometimes be too large and sometimes too small, making standard NP less consistently scaled than methods that account for the incoming activation norm.

Fan-in-scaled NP partially addresses this issue by accounting for the number of incoming units. If the incoming activations are independent, zero-mean, and unit-variance, then
\begin{equation}
\mathbb{E}\!\left[\left\|\mathbf{x}^{\ell-1}\right\|^2\right]
=
\mathbb{E}\!\left[\sum_{i=1}^{N_{\ell-1}} (x_i^{\ell-1})^2\right]
=
\sum_{i=1}^{N_{\ell-1}} \mathbb{E}\!\left[(x_i^{\ell-1})^2\right]
=
N_{\ell-1}.
\end{equation}
Thus, under these assumptions, the expected squared norm of the incoming activation vector grows with the fan-in. This is the same scaling argument used by \textcite{werfel2005learning}, who choose the NP perturbation variance to be a factor $N$ larger than the WP perturbation variance so that the two stochastic variants are directly comparable in their linear setting. If the incoming activations were explicitly normalized so that their squared norm matched the layer width, this correction would be exact up to the details of the normalization. However, this would introduce an additional architectural mechanism that is not part of the perturbation-learning methods reviewed here. Some reviewed work uses related stabilization techniques, such as weight normalization or input decorrelation, but these do not correspond to layer-wise activation normalization that fixes $\|\mathbf{x}^{\ell-1}\|^2$ in each forward pass. Moreover, common artificial normalization methods such as layer normalization typically include learnable scale and shift parameters, so they do not generally imply a fixed activation norm unless constrained to do so. Relying on such normalization would therefore make fan-in scaling depend on an extra assumption, and one whose biological interpretation is not established in the perturbation-learning literature considered here. In the setting studied in this thesis, activation norms are therefore allowed to vary across layers, samples, and training time. Fan-in scaling can therefore improve on standard NP by correcting for the typical size of a layer, but it still does not account for the actual incoming activation norm in a given forward pass.

IS-NP avoids this average correction by using the actual incoming activation norm in the current forward pass. Its perturbation scale is therefore adapted to the layer, input example, and current network state. This does not prove that IS-NP must outperform standard NP or fan-in-scaled NP in every practical setting, but it gives a clear scaling-based reason to expect more consistent behaviour when activation norms vary across layers, samples, or training time.
\subsection{Summary of theoretical contribution}

This chapter addressed the second research objective of the thesis: to further develop the theory of perturbation-based learning. Building on the gap identified in Chapter~2, we analysed the relationship between WP and NP at the estimator level in nonlinear multilayer networks. This led to the derivation of input-scaled node perturbation (IS-NP), a node perturbation rule that samples the pre-activation noise induced by Gaussian WP directly. We then showed that WP and IS-NP estimate the same perturbation-smoothed objective in expectation, and proved that IS-NP has variance smaller than or equal to WP for each coordinate. Finally, we compared IS-NP conceptually with standard NP and fan-in-scaled NP. This gave scaling-based reasons to expect IS-NP to behave more consistently, although this comparison was not a formal variance proof.

Chapters~2 and~3 therefore establish the theoretical foundation for the thesis. Chapter~2 addressed the first research objective by reviewing and clarifying the existing theory of perturbation-based learning, while Chapter~3 addressed the second objective. What remains is to address the third research objective: to evaluate how WP, standard NP, fan-in-scaled NP, and IS-NP perform in practice. The next chapter therefore turns from theory to experiments, where these four learning rules are compared empirically across the selected benchmark tasks.